\documentclass{article}
\usepackage[margin=2cm]{geometry}
\usepackage{amsmath}
\usepackage{graphicx}
\usepackage{booktabs}
\usepackage[utf8]{inputenc}
\usepackage{ragged2e}
\setlength{\emergencystretch}{3em}
\begin{document}
\sloppy

\noindent\hfill\includegraphics[height=1.8em]{assets/figure_header_logo.png}\par\noindent\hrule\vspace{4pt}

\subsection*{0.1Mass deformations}


In the following, we consider the coset with n = 1, i.e. a single vector multiplet. The coset
representative is expressed in terms of four scalars \$, i = 0, 1, 2, 3 via


\begin{equation}
L=\prod_{i=0}^{3}{c^{\phi}\dot{R}^{*}}\qquad (1)
\end{equation}


where K' are the non compact generators of SO(4,1); see for details. Note that \$o is an
SU(2)R singlet, while the other three scalars r form an SU(2)R triplet. The scalar potential
for this specific case can be obtained from and takes the following form


\begin{equation}
\begin{array}{c}{{V(\sigma,\dot{\phi}^{i})=-g^{2}e^{2\sigma}+\frac{1}{8}m e^{-6\sigma}\left[-32g e^{4\sigma}\cosh\phi^{0}\cosh\phi^{1}\cosh\phi^{2}\cosh\phi^{3}+8m\cosh^{2}\phi^{0}\right]}}\\ {{\ \ \ \ \ \ \ \ \left.+m\sinh^{2}\phi^{0}\left(-6+8\cosh^{2}\phi^{1}\cosh^{2}\phi^{2}\cosh^{2}\phi^{2}\cosh^{2}\phi^{2}\cosh(2\phi^{1}-\phi^{2})\right)+2\cosh(2\phi^{1})\right]}}\\ {{\cdot}}&{{\cdot\ \ \cdot\ \ \ .\ \ \ \ \ \ \ \ .\ \ \ \ \ \ \ \ \ \ \ \ \ \ .\ \ \ \ \ \ \ \ \ \ \ \ \ \ \ \ \ \ \ \ \ \ \ \ \ .\ \ \ \ \ \ \ \ \ \ \ \ \ \ \ \ \ \ \ \ \ \ \ \ \ \ \ \ \ \ \ \ \ \ \ \ \ \ \ \ \ \ \ \ \ \ \ \ \ \ \ \ \ \ \ \ \ \ \ \ \ \ \ \ \ \ \ \ \ \ \ \ \ \ \ \ \ \ \ \ \ \ \ \ \ \ \ \ \ \ \ \ \ \ \ \ \ \ \ \ \ \ \ \ \ \ \ \ \ \ \ \ \ \ \ \ \ \ \ \ \ \ \ \ \ \ \ \ \ \ \ \ \ \ \ \ \ \ \ \ \ \ \ \ \ \ \ \ \ \ \ \ \ \ \ \ \ \ \ \ \ \ \ \ \ \ \ \ \ \ \ \ \ \ \ \ \ \ \ \ \ \ \ \ \ \ \ \ \ \ \ \ \ \ \ \}}\end{array}
\end{equation}


The supersymmetric AdS6 vacuum is given by setting g = 3m and setting all scalars to.
vanish. The masses of the linearized scalar fluctuation around the AdS vacuum determine
the dimensions of the dual scalar operators in the SCFT via.


\begin{equation}
m^{2}l^{2}=\Delta(\Delta-5)~
\end{equation}


where l is the curvature radius of the AdS6 vacuum. For the scalars at hand, one finds


\begin{equation}
m_{\omega}^{2}l^{2}=-6\mathrm{\\\\\\\\\\\\\\\\\\\\}\ m_{\omega}^{2}l^{2}=-4\ \ \ \ \ \ \ \ \ \ \ \ \ \ \ m_{\omega}^{2}l^{2}=-6\ ,\ \ \ r=1,2,3\ \ \ \ \ \ \ \ \ \ \ \ \ \ \ (4)
\end{equation}


Hence the dimensions of the dual operators are.


\begin{equation}
\begin{array}{c c c}{{\Delta_{O_{\nu}}=3.\qquad\qquad}}&{{\qquad\Delta_{O_{\nu}}=4,\qquad\qquad}}&{{\Delta_{O_{\nu}}=3\ ,\ \ r=1,2,3\qquad\qquad}}&{{\qquad}}&{{\qquad(5)}}\end{array}
\end{equation}


In these CFT operators were expressed in terms of free hypermultiplets (i.e. the singleton
sector). The case of n = 1 corresponds to having a single free hypermultiplet, consisting of
four real scalars qA and two symplectic Majorana spinors '. Here I = 1, 2 is the SU(2)R
R-symmetry index and A = 1, 2 is the SU(2) flavor symmetry index. The gauge invariant
operators appearing in are related to these fundamental fields as follows,


\begin{equation}
\begin{array}{l l l l}{{{\cal O}_{\varphi}=(q^{*})^{A}{}_{t}q^{I}{}_{A},\quad}}&{{{\cal O}_{\varphi}=\bar{\psi}_{t}\psi^{I},\quad}}&{{{\cal O}_{\varphi^{*}}=(q^{*})^{A}{}_{t}(\sigma^{\prime})_{A}^{\;B}q^{I}{}_{B}\;,\;\;}}&{{r=1,2,3}}\\ {{-\pm\cdot\cdot\cdot{}_{-\cdot{}~c...}}}&{{-\cdot-\cdot-\cdot-\cdot-\cdot-\cdot-\cdot-\cdot-\cdot-\cdot-\cdot-\cdot-\cdot-\cdot-\cdot-\cdot-\cdot-\cdot-\cdot-\cdot-\cdot-\cdot-\cdot-\cdot-\cdot-\cdot-\cdot-\cdot-\cdot-\cdot-\cdot-\cdot-\cdot-\cdot-\cdot-\cdot-\cdot-\cdot-\cdot-\cdot-\cdot-\cdot-\cdot-\cdot-\cdot-\cdot-\cdot-\cdot-\cdot-\cdot-\cdot-\cdot-\cdot-\cdot-\cdot-\cdot-\cdot-\cdot-\cdot-\ldots}}}&{\leq-\leq\leq\leq\ldots}&{\leq-\leq-\leq\ldots.}&{\leq\ldots\ldots-\ldots\ldots\ldots\ldots\ldots\ldots\ldots\ldots-\ldots-\ldots\ldots\ldots\ldots-\ldots-\ldots\ldots\ldots\ldots-\ldots\ldots\ldots\ldots\ldots\ldots\leq-\ldots\ldots\end{array}
\end{equation}


Note that the first two operators correspond to mass terms for the scalars and fermions

\end{document}
